%!TeX root=../princesstop.tex
\chapter{Lottie}

\begin{figure}[t!]
\centering
\includegraphics[width=\linewidth]{filigree}
\end{figure}

\lettrine[lines=5]{I}{f} Sara had been a different kind of child, the life she led at Miss Minchin's Select Seminary for the next few years would not have been at all good for her. She was treated more as if she were a distinguished guest at the establishment than as if she were a mere little girl. If she had been a self-opinionated, domineering child, she might have become disagreeable enough to be unbearable through being so much indulged and flattered. If she had been an indolent child, she would have learned nothing. Privately Miss Minchin disliked her, but she was far too worldly a woman to do or say anything which might make such a desirable pupil wish to leave her school. She knew quite well that if Sara wrote to her papa to tell him she was uncomfortable or unhappy, Captain Crewe would remove her at once. Miss Minchin's opinion was that if a child were continually praised and never forbidden to do what she liked, she would be sure to be fond of the place where she was so treated. Accordingly, Sara was praised for her quickness at her lessons, for her good manners, for her amiability to her fellow pupils, for her generosity if she gave sixpence to a beggar out of her full little purse; the simplest thing she did was treated as if it were a virtue, and if she had not had a disposition and a clever little brain, she might have been a very self-satisfied young person. But the clever little brain told her a great many sensible and true things about herself and her circumstances, and now and then she talked these things over to Ermengarde as time went on.

»Things happen to people by accident,« she used to say. »A lot of nice accidents have happened to me. It just \textsc{happened} that I always liked lessons and books, and could remember things when I learned them. It just happened that I was born with a father who was beautiful and nice and clever, and could give me everything I liked. Perhaps I have not really a good temper at all, but if you have everything you want and everyone is kind to you, how can you help but be good-tempered? I don't know«—looking quite serious—»how I shall ever find out whether I am really a nice child or a horrid one. Perhaps I'm a \textsc{hideous} child, and no one will ever know, just because I never have any trials.«

»Lavinia has no trials,« said Ermengarde, stolidly, »and she is horrid enough.«

Sara rubbed the end of her little nose reflectively, as she thought the matter over.

»Well,« she said at last, »perhaps—perhaps that is because Lavinia is\textsc{growing}.« This was the result of a charitable recollection of having heard Miss Amelia say that Lavinia was growing so fast that she believed it affected her health and temper.

Lavinia, in fact, was spiteful. She was inordinately jealous of Sara. Until the new pupil's arrival, she had felt herself the leader in the school. She had led because she was capable of making herself extremely disagreeable if the others did not follow her. She domineered over the little children, and assumed grand airs with those big enough to be her companions. She was rather pretty, and had been the best-dressed pupil in the procession when the Select Seminary walked out two by two, until Sara's velvet coats and sable muffs appeared, combined with drooping ostrich feathers, and were led by Miss Minchin at the head of the line. This, at the beginning, had been bitter enough; but as time went on it became apparent that Sara was a leader, too, and not because she could make herself disagreeable, but because she never did.

»There's one thing about Sara Crewe,« Jessie had enraged her »best friend« by saying honestly, »she's never »grand« about herself the least bit, and you know she might be, Lavvie. I believe I couldn't help being—just a little—if I had so many fine things and was made such a fuss over. It's disgusting, the way Miss Minchin shows her off when parents come.«

»»Dear Sara must come into the drawing room and talk to Mrs Musgrave about India,«« mimicked Lavinia, in her most highly flavoured imitation of Miss Minchin. »»Dear Sara must speak French to Lady Pitkin. Her accent is so perfect.« She didn't learn her French at the Seminary, at any rate. And there's nothing so clever in her knowing it. She says herself she didn't learn it at all. She just picked it up, because she always heard her papa speak it. And, as to her papa, there is nothing so grand in being an Indian officer.«

»Well,« said Jessie, slowly, »he's killed tigers. He killed the one in the skin Sara has in her room. That's why she likes it so. She lies on it and strokes its head, and talks to it as if it was a cat.«

»She's always doing something silly,« snapped Lavinia. »My mamma says that way of hers of pretending things is silly. She says she will grow up eccentric.«

It was quite true that Sara was never »grand.« She was a friendly little soul, and shared her privileges and belongings with a free hand. The little ones, who were accustomed to being disdained and ordered out of the way by mature ladies aged ten and twelve, were never made to cry by this most envied of them all. She was a motherly young person, and when people fell down and scraped their knees, she ran and helped them up and patted them, or found in her pocket a bonbon or some other article of a soothing nature. She never pushed them out of her way or alluded to their years as a humiliation and a blot upon their small characters.

»If you are four you are four,« she said severely to Lavinia on an occasion of her having—it must be confessed—slapped Lottie and called her »a brat;« »but you will be five next year, and six the year after that. And,« opening large, convicting eyes, »it takes sixteen years to make you twenty.«

»Dear me,« said Lavinia, »how we can calculate!« In fact, it was not to be denied that sixteen and four made twenty—and twenty was an age the most daring were scarcely bold enough to dream of.

So the younger children adored Sara. More than once she had been known to have a tea party, made up of these despised ones, in her own room. And Emily had been played with, and Emily's own tea service used—the one with cups which held quite a lot of much-sweetened weak tea and had blue flowers on them. No one had seen such a very real doll's tea set before. From that afternoon Sara was regarded as a goddess and a queen by the entire alphabet class.

Lottie Legh worshipped her to such an extent that if Sara had not been a motherly person, she would have found her tiresome. Lottie had been sent to school by a rather flighty young papa who could not imagine what else to do with her. Her young mother had died, and as the child had been treated like a favourite doll or a very spoiled pet monkey or lap dog ever since the first hour of her life, she was a very appalling little creature. When she wanted anything or did not want anything she wept and howled; and, as she always wanted the things she could not have, and did not want the things that were best for her, her shrill little voice was usually to be heard uplifted in wails in one part of the house or another.

Her strongest weapon was that in some mysterious way she had found out that a very small girl who had lost her mother was a person who ought to be pitied and made much of. She had probably heard some grown-up people talking her over in the early days, after her mother's death. So it became her habit to make great use of this knowledge.

The first time Sara took her in charge was one morning when, on passing a sitting room, she heard both Miss Minchin and Miss Amelia trying to suppress the angry wails of some child who, evidently, refused to be silenced. She refused so strenuously indeed that Miss Minchin was obliged to almost shout—in a stately and severe manner—to make herself heard.

»What \textsc{is} she crying for?« she almost yelled.

»Oh—oh—oh!« Sara heard; »I haven't got any mam—ma-a!«

»Oh, Lottie!« screamed Miss Amelia. »Do stop, darling! Don't cry! Please don't!«

»Oh! Oh! Oh! Oh! Oh!« Lottie howled tempestuously. »Haven't—got—any—mam—ma-a!«

»She ought to be whipped,« Miss Minchin proclaimed. »You \textsc{shall} be whipped, you naughty child!«

Lottie wailed more loudly than ever. Miss Amelia began to cry. Miss Minchin's voice rose until it almost thundered, then suddenly she sprang up from her chair in impotent indignation and flounced out of the room, leaving Miss Amelia to arrange the matter.

Sara had paused in the hall, wondering if she ought to go into the room, because she had recently begun a friendly acquaintance with Lottie and might be able to quiet her. When Miss Minchin came out and saw her, she looked rather annoyed. She realized that her voice, as heard from inside the room, could not have sounded either dignified or amiable.

»Oh, Sara!« she exclaimed, endeavouring to produce a suitable smile.

»I stopped,« explained Sara, »because I knew it was Lottie—and I thought, perhaps—just perhaps, I could make her be quiet. May I try, Miss Minchin?«

»If you can, you are a clever child,« answered Miss Minchin, drawing in her mouth sharply. Then, seeing that Sara looked slightly chilled by her asperity, she changed her manner. »But you are clever in everything,« she said in her approving way. »I dare say you can manage her. Go in.« And she left her.

When Sara entered the room, Lottie was lying upon the floor, screaming and kicking her small fat legs violently, and Miss Amelia was bending over her in consternation and despair, looking quite red and damp with heat. Lottie had always found, when in her own nursery at home, that kicking and screaming would always be quieted by any means she insisted on. Poor plump Miss Amelia was trying first one method, and then another.

»Poor darling,« she said one moment, »I know you haven't any mamma, poor\longdash« Then in quite another tone, »If you don't stop, Lottie, I will shake you. Poor little angel! There—! You wicked, bad, detestable child, I will smack you! I will!«

Sara went to them quietly. She did not know at all what she was going to do, but she had a vague inward conviction that it would be better not to say such different kinds of things quite so helplessly and excitedly.

»Miss Amelia,« she said in a low voice, »Miss Minchin says I may try to make her stop—may I?«

Miss Amelia turned and looked at her hopelessly. »Oh, \textsc{do} you think you can?« she gasped.

»I don't know whether I \textsc{can}«, answered Sara, still in her half-whisper; »but I will try.«

Miss Amelia stumbled up from her knees with a heavy sigh, and Lottie's fat little legs kicked as hard as ever.

»If you will steal out of the room,« said Sara, »I will stay with her.«

»Oh, Sara!« almost whimpered Miss Amelia. »We never had such a dreadful child before. I don't believe we can keep her.«

But she crept out of the room, and was very much relieved to find an excuse for doing it.

Sara stood by the howling furious child for a few moments, and looked down at her without saying anything. Then she sat down flat on the floor beside her and waited. Except for Lottie's angry screams, the room was quite quiet. This was a new state of affairs for little Miss Legh, who was accustomed, when she screamed, to hear other people protest and implore and command and coax by turns. To lie and kick and shriek, and find the only person near you not seeming to mind in the least, attracted her attention. She opened her tight-shut streaming eyes to see who this person was. And it was only another little girl. But it was the one who owned Emily and all the nice things. And she was looking at her steadily and as if she was merely thinking. Having paused for a few seconds to find this out, Lottie thought she must begin again, but the quiet of the room and of Sara's odd, interested face made her first howl rather half-hearted.

»I—haven't—any—ma—ma—ma-a!« she announced; but her voice was not so strong.

Sara looked at her still more steadily, but with a sort of understanding in her eyes.

»Neither have I,« she said.

This was so unexpected that it was astounding. Lottie actually dropped her legs, gave a wriggle, and lay and stared. A new idea will stop a crying child when nothing else will. Also it was true that while Lottie disliked Miss Minchin, who was cross, and Miss Amelia, who was foolishly indulgent, she rather liked Sara, little as she knew her. She did not want to give up her grievance, but her thoughts were distracted from it, so she wriggled again, and, after a sulky sob, said, »Where is she?«

Sara paused a moment. Because she had been told that her mamma was in heaven, she had thought a great deal about the matter, and her thoughts had not been quite like those of other people.

»She went to heaven,« she said. »But I am sure she comes out sometimes to see me—though I don't see her. So does yours. Perhaps they can both see us now. Perhaps they are both in this room.«

Lottie sat bolt upright, and looked about her. She was a pretty, little, curly-headed creature, and her round eyes were like wet forget-me-nots. If her mamma had seen her during the last half-hour, she might not have thought her the kind of child who ought to be related to an angel.

Sara went on talking. Perhaps some people might think that what she said was rather like a fairy story, but it was all so real to her own imagination that Lottie began to listen in spite of herself. She had been told that her mamma had wings and a crown, and she had been shown pictures of ladies in beautiful white nightgowns, who were said to be angels. But Sara seemed to be telling a real story about a lovely country where real people were.

»There are fields and fields of flowers,« she said, forgetting herself, as usual, when she began, and talking rather as if she were in a dream, »fields and fields of lilies—and when the soft wind blows over them it wafts the scent of them into the air—and everybody always breathes it, because the soft wind is always blowing. And little children run about in the lily fields and gather armfuls of them, and laugh and make little wreaths. And the streets are shining. And people are never tired, however far they walk. They can float anywhere they like. And there are walls made of pearl and gold all round the city, but they are low enough for the people to go and lean on them, and look down onto the earth and smile, and send beautiful messages.«

Whatsoever story she had begun to tell, Lottie would, no doubt, have stopped crying, and been fascinated into listening; but there was no denying that this story was prettier than most others. She dragged herself close to Sara, and drank in every word until the end came—far too soon. When it did come, she was so sorry that she put up her lip ominously.

»I want to go there,« she cried. »I—haven't any mamma in this school.«

Sara saw the danger signal, and came out of her dream. She took hold of the chubby hand and pulled her close to her side with a coaxing little laugh.

»I will be your mamma,« she said. »We will play that you are my little girl. And Emily shall be your sister.«

Lottie's dimples all began to show themselves.

»Shall she?« she said.

»Yes,« answered Sara, jumping to her feet. »Let us go and tell her. And then I will wash your face and brush your hair.«

To which Lottie agreed quite cheerfully, and trotted out of the room and upstairs with her, without seeming even to remember that the whole of the last hour's tragedy had been caused by the fact that she had refused to be washed and brushed for lunch and Miss Minchin had been called in to use her majestic authority.

And from that time Sara was an adopted mother.

\begin{a4}
	\enlargethispage{\baselineskip}
\end{a4}